\documentclass[10pt,a4paper]{article}

\usepackage[margin=1.05in]{geometry}
\usepackage{setspace}
\usepackage{amsmath,amssymb,amsthm}
\usepackage{graphicx}
\usepackage{enumitem}
\usepackage{titlesec}
\usepackage[expansion=false]{microtype}
\usepackage{fancyhdr}
\usepackage{hyperref}
\hypersetup{colorlinks=true,linkcolor=black,urlcolor=black,citecolor=black}

\graphicspath{{../set/}}

\titleformat{\section}{\normalsize\bfseries}{}{0em}{}
\titleformat{\subsection}{\small\bfseries}{}{0em}{}
\titlespacing*{\section}{0pt}{1.35em}{0.5em}
\titlespacing*{\subsection}{0pt}{0.9em}{0.3em}

\theoremstyle{plain}
\newtheorem{proposition}{Proposition}
\theoremstyle{definition}
\newtheorem{definition}{Definition}
\theoremstyle{remark}
\newtheorem{observation}{Observation}

\setstretch{1.14}
\setlist[itemize]{leftmargin=1.15em,itemsep=0.1em,topsep=0.2em}

\pagestyle{fancy}
\fancyhf{}
\fancyhead[L]{\small\itshape Elements of Position}
\fancyhead[R]{\small\itshape Without an Angle}
\fancyfoot[C]{\small\thepage}
\renewcommand{\headrulewidth}{0pt}

\begin{document}

\begin{center}
\vspace*{2.0em}
{\large Elements of Position}\\[1.6em]
{\LARGE\bfseries Without an Angle}\\[0.6em]
{\small Two moves, the pair they force, and where its two readings meet}\\[1.8em]
{\normalsize Justin Erholtz}\\[0.35em]
{\small 19 September 2026}
\end{center}

\vspace{2.0em}

\begin{center}
{\large The operators do not contain the joint.}
\end{center}

\vspace{1.6em}

% ======================================================================
\section*{Before the proofs}
\addcontentsline{toc}{section}{Before the proofs}
% ======================================================================

Every later part of this series opens quickly and stays inside its proofs.
This one does not, on purpose, because this is the part the others are built on, and a foundation is worth a longer look before anything is set on top of it.

I want to say plainly how this object got made, because the how is not separate from what can be trusted in the what.
It did not start as a proof.
It started as two lines of arithmetic, run a few hundred thousand times, and a rule I tried to hold to throughout: a number does not get to stand until it has been produced twice, by two different routes, and the two routes agree.
A claim proved on paper and a claim read off a table are different kinds of claim, and the difference is not a formality.
It is the whole discipline this series calls object, sheet, and entrance, and it was earned the slow way, by being wrong about the boundary between them more than once and fixing it each time rather than smoothing over it.

Two of those corrections are worth naming, because they shaped what this part is willing to say and how carefully it says it.

The first was a habit of language, not of mathematics.
It is easy, faced with an unfamiliar construction, to reach for one of two lazy verdicts: that it must already be known, or that it must be new.
Both are a way of not looking closely.
The honest version takes longer: check what the operators alone force, check it again by hand, and only then ask what it happens to be called elsewhere.
Every identity in this part was checked that way, in that order.
Where the order was reversed even once --- where a name got reached for before the check was finished --- it was a mistake, and it is not repeated here.

The second correction came from depth rather than direction.
The leftover \(\tau(m)\) has a closed-form sum (Proposition~\ref{prop:tausum} below), and at the scale of the sheets in this part it agrees with that closed form to every digit shown.
Run the same walk far past that scale --- hundreds of times deeper, into the millions of hits --- and a naive running total of \(\tau(m)\) quietly disagrees with the \emph{integral} \(\int_1^N\tau\) by an amount that grows with the depth of the run.
The first reading of that gap, kept here only long enough to correct, was that the addition had failed: summing millions of terms into a number many orders of magnitude larger than any one of them loses precision on every step.
It had not.
A Kahan-compensated running total --- carried with a separate correction term meant to catch exactly that kind of loss --- disagrees from the integral by the same amount as the naive total, to every digit tested.
Nothing was lost in the addition.
What was missing was a term in the comparison: the integral is only the leading part of the discrete sum, and the gap between them is not noise but the Euler--Maclaurin correction Proposition~\ref{prop:wedge} names below, a quantity as fixed as \(\tau\) itself and one this series had simply not yet written down.
Once the trapezoidal correction is restored the residual stops growing with \(N\) and settles on \(\gamma/2-\tfrac14\), the constant Proposition~\ref{prop:wedge} derives.
The lesson survives the correction, sharpened rather than softened: a sum this series reports is always the discrete closed form of Proposition~\ref{prop:tausum}, checked against a direct count at a scale small enough for the count itself to be trustworthy --- and never a running total left to grow unchecked against the wrong comparison, which will disagree with almost anything, precision intact.

With that said, the second thing worth doing plainly, before any proof below, is placing this object next to the mathematics it touches, rather than letting a reader wonder.

Almost nothing proved in this part is new.
The harmonic conjugate of Proposition~\ref{prop:harmonic} is classical projective geometry: any involution of a line with two fixed points sends every point to a partner in a fixed cross-ratio of \(-1\) with those two points, a fact that predates this series by roughly two thousand years in spirit and about a century and a half in the specific vocabulary of cross-ratio itself \cite{coxeter1969}.
The two-integer gap structure of \S\ref{sec:twoint} is a special case of a fact conjectured by Steinhaus and proved independently, within a few years of one another, by several mathematicians studying the fractional parts of multiples of an irrational number \cite{sos1958,swierczkowski1958}.
Building a rotation from a shear and a renormalization, rather than calling sine and cosine directly, is a known technique with its own name and its own lineage in digital computing, developed for a navigation computer, for entirely different reasons, decades before this series existed \cite{volder1959}.
Summing a slowly growing quantity and comparing the sum with a smooth integral is Euler and Maclaurin's, older than the notation used here to state it \cite{apostol1976}.
None of that is a discovery, and this series has said so on its first page.

What is offered instead is the object: one generator, refusing to name an angle anywhere inside its own two moves, that produces all four of the results above as consequences of a single walk, measured against one another rather than assembled to order.
Whether that assembly is worth the years this series intends to give it is a claim the reader can weigh against the tables and the proofs that follow --- not against this paragraph.

\vspace{1.2em}
\noindent
\textit{A note.}
Claims of two kinds are kept apart throughout.
A claim proved from the two moves and the count is an identity, true for every admissible kick and every start on the diameter.
A claim that is a count --- a gap, a word, a mean taken over one run --- is reported from a table, and the table is named.
The recorded run cited throughout this part is \(h=0.002\), \(200\,000\) steps, \(127\) crossings.
That run is an example.
A different kick, or a longer run, writes a new table.
It does not revise an identity.
No angle enters the step.
A length that appears only after a table exists may be named as a comparison, and not before.

\vspace{0.7em}
\noindent
\textit{A note on names used here.}
Shear \(T_h\), bind \(N\), step --- the two moves and their composition.
Hit --- a crossing of the diameter the walk began on; the count of hits is \(m\).
Pair --- \(r=m\), \(1/r=1/m\), written on that count.
Lock --- the hit \(m=1\), where the pair meets itself.
Leftover \(\tau\) --- the gap between the arithmetic mean and the geometric mean of a pair.

\newpage
\tableofcontents
\newpage

% ======================================================================
\section{What is in the operators}
% ======================================================================

Start on the diameter: \((1,0)\).
A heading is a pair of coordinates whose length, after the bind, is \(1\).

\begin{definition}[Shear]
For a kick \(h>0\),
\[
T_h(x,y)=(x-hy,\; y+hx).
\]
\end{definition}

\begin{definition}[Bind]
\[
N(x,y)=\frac{(x,y)}{\sqrt{x^2+y^2}}.
\]
\end{definition}

\begin{definition}[Step]
\[
\mathrm{step}=N\circ T_h.
\]
\end{definition}

No angle is an argument of \(T_h\) or of \(N\).
The shear is an equal, opposite mix of the two coordinates.
The bind is length.
Read separately, neither move says anything about a turn.
Composed and repeated they produce a heading that comes back to the line it started on.
That return is counted.
It is not named as a turn inside the step --- though it is worth saying, once, that the composition is exactly a rotation by the fixed angle \(\arctan h\), a shear-and-renormalize construction with its own history in the fast computation of trigonometric functions rather than in geometry \cite{volder1959}.
That fact is never used inside the walk.
It is used only to name a comparison, later, once a table exists to compare against.

If the last heading already had length \(1\), the unbound length after the shear is fixed by \(h\) alone:
\[
\bigl|T_h(u)\bigr|=\sqrt{1+h^2}.
\]
The leftover of one kick is therefore \(\sqrt{1+h^2}-1\), constant in the step index.
On the recorded run this is \(1.999998\times 10^{-6}\).
The Euclidean distance between two successive bound headings is likewise a function of \(h\) only.
On the recorded run it is \(0.001999997\).
Those two numbers describe the grain of one run.
They are not a position on the pair.

\begin{observation}
An unequal mix \((x-ay,\, y+bx)\) still returns to the line, but mean gap times kick tracks a different comparison.
A one-sided mix from this start does not return.
The equal opposite shear is a choice of kind.
It is a right angle in the mix.
It is not a finished turn in the mix.
\end{observation}

\begin{figure}[htbp]
\centering
\includegraphics[width=0.92\textwidth]{W1_heading.png}
\caption{Sheet W1. The bind, the cut, and the stations the kernel recorded. Station at $(1,0)$, $h=0.002$. Hits lie on $y=0$.}
\label{fig:plan}
\end{figure}

% ======================================================================
\section{The hit}
% ======================================================================

\begin{definition}[Hit]
A heading starting at \((1,0)\) is walked by repeated steps.
A hit is a step at which the heading's second coordinate changes sign.
\end{definition}

\begin{definition}[First gap]
\(H_0\) is the step index of the first hit.
\end{definition}

On the recorded run, \(H_0=1571\), side down, heading
\[
(-0.99999992,\; -0.00040316).
\]
The heading is on the far side of the line, one kick of vertical off it.

Stand on the line you count and the first event is the far side of that line.
Stand off the line and the first gap is only how far the start was from the wall.
This run stands on the diameter.

\begin{observation}
The distance from the start heading to the first hit heading is \(1.99999996\).
The distance from the first hit heading to the second is the same figure.
A heading exactly on the far side of the axis would give exactly \(2\).
The walk lands one kick of vertical short of that, every time, because a hit is a sign change and not a landing.
\end{observation}

% ======================================================================
\section{The pair and the lock}
\label{sec:pairlock}
% ======================================================================

After the first hit the scale of a hit is its index \(m=1,2,3,\ldots\).

\begin{definition}[Pair]
\[
r=m,\qquad \frac1r=\frac1m.
\]
\end{definition}

The pair is written on the integer \(m\).
It does not read the heading.
Once \(m\) is named, the walk is no longer required.

\begin{proposition}
\label{prop:pair}
For every \(m\ge 1\),
\[
r\cdot\frac1r=1,\qquad
\mathrm{GM}\!\left(r,\frac1r\right)=1.
\]
\end{proposition}

\begin{proof}
The product is immediate.
The geometric mean is \(\sqrt{r\cdot(1/r)}=\sqrt{1}=1\).
\end{proof}

\begin{definition}[Leftover]
\[
\tau(m)=\mathrm{AM}-\mathrm{GM}
=\frac{m+1/m}{2}-1.
\]
\end{definition}

\begin{proposition}
\label{prop:tau}
\[
\tau(m)=\frac{(m-1)^2}{2m}.
\]
\(\tau(m)=0\) if and only if \(m=1\).
\end{proposition}

\begin{proof}
\[
\frac{m+1/m}{2}-1
=\frac{m^2+1-2m}{2m}
=\frac{(m-1)^2}{2m}.
\]
The numerator vanishes only at \(m=1\).
\end{proof}

\begin{proposition}
\label{prop:tau-inv}
\(\tau(m)=\tau(1/m)\) for every \(m>0\).
\end{proposition}

\begin{proof}
Substitute.
\end{proof}

Hit \(m=1\) is the lock.
The pair is \((1,1)\).
The leftover is \(0\).
The count has moved once and has not yet had room to differ from its own inverse.

Hit \(m=2\) is the first pair that is two places: \((2,1/2)\), arithmetic mean \(5/4\), leftover \(1/4\).
Both ends now exist.

Nothing after the lock is a new kind of figure.
Every later \(m\) writes the same pair, the same product, the same leftover formula, at a larger \(r\).

\begin{figure}[htbp]
\centering
\includegraphics[width=0.92\textwidth]{D1_lock.png}
\caption{Sheet D1. The lock: pair $(1,1)$ at $m=1$. The bind is the pair.}
\label{fig:lock}
\end{figure}

\begin{figure}[htbp]
\centering
\includegraphics[width=0.92\textwidth]{D2_pair.png}
\caption{Sheet D2. The first two-place pair $(2,1/2)$ at $m=2$, on the home wall.}
\label{fig:pair}
\end{figure}

% ======================================================================
\section{The outward point}
\label{sec:outward}
% ======================================================================

At every hit the pair is not only two numbers.
It is a point.
\[
P_m = m\,u_m,
\]
where \(u_m\) is the bound heading at that hit.
\(P_1=u_1\) is the lock itself.
Every later \(P_m\) is the same construction, walked out along a different ray, at a distance the count itself supplies.

\begin{figure}[htbp]
\centering
\includegraphics[width=0.92\textwidth]{P1_plan.png}
\caption{Sheet P1. Plan of the model looking down the count. Outer locus $P_m=m u_m$ in ink; inverse locus $Q_m=u_m/m$ in the accent; bind of radius $1$; cut $y=0$.}
\label{fig:modelplan}
\end{figure}

\subsection{Four points on the line}
\label{sec:fourpoints}

At a hit the heading lies on the cut \(y=0\) to the residual of one kick.
Four points sit on that line, written as signed abscissae along the ray: the inner mark \(1/m\), the outer mark \(m\), the unit mark \(+1\), the opposite pole \(-1\).

\begin{proposition}
\label{prop:harmonic}
For every \(m>1\),
\[
\left(\frac1m,\, m;\, 1,\, -1\right)=-1.
\]
\end{proposition}

\begin{proof}
\[
\frac{(1-1/m)/(1-m)}{(-1-1/m)/(-1-m)}
=\frac{(m-1)/m}{1-m}\cdot\frac{-(m+1)}{-(m+1)/m}
=-1.
\]
\end{proof}

On every live row of the recorded run the value is \(-1\) to machine precision.
The four points are the pair and the two marks the diameter already fixed.
This is not a coincidence of this particular pair: any involution of a line with two fixed points sends every point to its partner in a constant cross-ratio of \(-1\) with those two fixed points \cite{coxeter1969}.
The map \(z\mapsto 1/z\) is exactly such an involution, fixing \(1\) and \(-1\); Proposition~\ref{prop:harmonic} is the instance the walk happens to write, not an exception to the rule.

A neighboring four on the same line,
\[
\bigl(m-1,\, m+1;\, m,\, 0\bigr)=\frac{m+1}{m-1},
\]
depends on \(m\), and does not tend to a fixed value as an involution's does; it grows toward \(1\) only in the sense of approaching it from above as \(m\to\infty\).
Anchor both reference points on a tick instead of the origin, and the dependence sharpens rather than fades:

\begin{proposition}
\label{prop:tickorigin}
For every \(m>2\),
\[
\bigl(m-1,\,m+1;\,m,\,1\bigr)=\frac{m}{2-m}.
\]
\end{proposition}

\begin{proof}
Direct evaluation of the cross-ratio formula; the \(m\) does not cancel because one reference point, \(1\), is fixed while the other, \(m\), moves with the count.
\end{proof}

The value tends to \(-1\) as \(m\to\infty\) --- the same limit Proposition~\ref{prop:harmonic} reaches exactly at every finite \(m\) --- but only ever approaches it here, never reaches it at any finite hit.
The lock against the pair is constant.
Ticks against a fixed point that is not the pair's own fixed point are not, and the rate at which they fail to be constant is exactly \(m/(2-m)\).

\subsection{Ticks on the generator}
\label{sec:generator}

A different four, built entirely from evenly spaced integers rather than from the pair, is constant for a different reason.

\begin{proposition}
\label{prop:ladder}
For integers \(a\ne b\) and \(m>\max(a,b)\),
\[
\bigl(m-a,\,m+a;\,m-b,\,m+b\bigr)=\left(\frac{b-a}{a+b}\right)^{2}.
\]
\end{proposition}

\begin{proof}
Direct evaluation; the \(m\) cancels in the four-term formula before the square, since all four points shift by the same \(m\).
\end{proof}

In particular,
\[
\bigl(m-2,\,m+2;\,m-1,\,m+1\bigr)=\frac19,
\]
checked exactly at \(m=3,4,10,50,1000\).
This is a different invariance than Proposition~\ref{prop:harmonic}'s.
The pair is constant because \(1\) and \(-1\) are the fixed points of the same map that sends \(m\) to \(1/m\).
The generator ladder is constant because shifting four evenly spaced points by a common amount is itself an isometry of the line, and cross-ratio does not see translations.
Two unrelated reasons produce the same word, ``constant,'' on two unrelated families of four points; Proposition~\ref{prop:tickorigin} is the reminder that a four which is neither an involution's orbit nor an evenly spaced ladder has no such protection.

\begin{figure}[htbp]
\centering
\includegraphics[width=0.92\textwidth]{D5_four_points.png}
\caption{Sheet D5. The pair's harmonic four ($1/m$, $m$; $1$, $-1$) beside the generator's evenly spaced four ($m-2$, $m+2$; $m-1$, $m+1$), drawn on parallel rays at $m=10$.}
\label{fig:fourpoints}
\end{figure}

\subsection{The count pays with two integers}
\label{sec:twoint}

After the first hit the recorded run writes only two gap values: \(1570\) and \(1571\).
On \(127\) hits, \(1571\) occurs \(102\) times and \(1570\) occurs \(25\) times.
Mean gap \(1570.80315\).
Mean gap times the kick: \(3.14160630\).

The walk cannot land on one length, so it writes both neighbors.
This is the discrete case of a theorem conjectured by Steinhaus and proved independently by several authors studying the distribution of \(\{n\alpha\}\) for irrational \(\alpha\): the return times of a fixed rotation to an interval take at most two values, adjacent integers determined by the rotation's own angle \cite{sos1958,swierczkowski1958}.
The walk here is that rotation, and the interval is a half-turn.

A named comparison after the table exists:
\[
\pi=3.14159265\ldots,\qquad
\lambda:=\frac{\pi}{\arctan h}=1570.798421.
\]
Those names sit between \(1570\) and \(1571\).
\(H_0\cdot h=3.142\) is the first gap only --- one long step, coarser than the mean.

On this run \(\{\lambda\}=0.798421>1/2\), so the long letter dominates.
The remaining word of long and short gaps is twenty-five copies of LLLLS, then LL, because the step budget ended mid-word.
That five is not a second figure invented to describe the word; it is the denominator of the best simple approximation to \(1-\{\lambda\}=0.201579\), whose continued fraction begins \([0;4,1,24,\ldots]\) with second convergent \(1/5\).
A short gap fires, on average, once every five hits, because \(1/5\) is the closest simple fraction to the short gap's own frequency.
A different \(h\) writes a different fraction and a different period, and the period is computable from \(h\) before the walk is ever run.

\begin{observation}
\label{obs:exactclosure}
Choose \(h=\tan(\pi/n)\) for an integer \(n\).
Then \(\lambda=n\) exactly, and every gap after the first hit equals \(n\), with no second letter, checked to \(200{,}000\) hits at \(n=100\).
\end{observation}

The two-letter alphabet of this section is what the walk writes at almost every kick.
It is not forced at every kick.
At the exact mixes of Observation~\ref{obs:exactclosure} the walk closes on a single integer, the same way a rotation by a rational multiple of \(\pi\) closes a finite orbit instead of an infinite one --- a measure-zero set of kicks, but not an empty one.

\begin{figure}[htbp]
\centering
\includegraphics[width=0.92\textwidth]{SCH1_integers.png}
\caption{Sheet SCH1. The two-integer schedule of the recorded run.}
\label{fig:sch1}
\end{figure}

\subsection{The heading ratchets and resets}
\label{sec:ratchet}

On a long hit the heading is off the axis by about a kick, then two, then three, then four.
On the next short hit it is put back.

\begin{observation}
On every short hit of the recorded run,
\[
\frac{|h_y|}{m}=\alpha,\qquad
\alpha\approx 3.15762\times 10^{-6},
\]
to the digits of the mean, from \(m=5\) through \(m=125\).
\end{observation}

Snapped to the event, the heading is always a horizon hit, by definition.
Held against the first gap as a ruler, it walks off that ruler until a short gap settles the account.
The two readings are two rulers on the same walk.

Call the gap between an exact landing on the axis and the heading the operators actually deliver the \emph{entrance residual}.
It shrinks as \(h\) shrinks and never reaches zero.

\begin{figure}[htbp]
\centering
\includegraphics[width=0.92\textwidth]{E1_elevation.png}
\caption{Sheet E1. Elevation: $m\cdot h_y$ against $m$, showing the ratchet on long hits and the reset toward the cut on short hits (accent).}
\label{fig:elevation}
\end{figure}

\begin{figure}[htbp]
\centering
\includegraphics[width=0.92\textwidth]{D4_cycle.png}
\caption{Sheet D4. One LLLLS cycle in $m\cdot h_y$, hits $m=1$ through $m=5$.}
\label{fig:cycle}
\end{figure}

\subsection{Chord and path}
\label{sec:chordpath}

\begin{proposition}
\label{prop:chord}
If the heading at hit \(m-1\) is exactly opposite the heading at hit \(m\), then
\[
|P_m-P_{m-1}|=2m-1.
\]
\end{proposition}

\begin{proof}
Opposite unit headings \(u\) and \(-u\) give
\[
\bigl|m(-u)-(m-1)u\bigr|=(m+m-1)|u|=2m-1.
\]
\end{proof}

The hypothesis is never exactly true, by the residual already named.
On the recorded run the chord at \(m=2\) is \(3.000000\), at \(m=3\) it is \(5.000000\), at \(m=4\) it is \(7.000000\).
The short-gap rows differ from \(2m-1\) by a few parts in \(10^6\), the same size as \(\alpha\).

\begin{proposition}[Radius cancels]
\label{prop:cancel}
Let the chord above be drawn across a half-turn, and let the path be the first-order displacement of a radius growing linearly across that half-turn, at rate \(h\) per step.
Then
\[
\frac{\mathrm{chord}}{\mathrm{path}}=\frac{2}{(\mathrm{gap})\cdot h}.
\]
\end{proposition}

The radius cancels.
The ratio does not depend on which wall is being read.
On the recorded run the two bands are \(2/(1570h)=0.63694\) and \(2/(1571h)=0.63654\).
A named comparison, \(2/\pi=0.63662\ldots\), sits between them.
The last row of this run is a long hit, so the last ratio is the lower band, not a leftover angle treated as a second constant.

The path used here is first-order.
It is not a summed displacement of \(P\).

\begin{figure}[htbp]
\centering
\includegraphics[width=0.92\textwidth]{SCH2_bands.png}
\caption{Sheet SCH2. Chord over first-order path, two bands, with $2/\pi$ between them.}
\label{fig:bands}
\end{figure}

\subsection{The leftover, summed}
\label{sec:tausum}

\begin{proposition}
\label{prop:tausum}
\[
\sum_{m=1}^{N}\tau(m)=\frac{N(N+1)}{4}+\frac{H_N}{2}-N,
\]
where \(H_N=\sum_{m=1}^{N}1/m\).
\end{proposition}

\begin{proof}
Sum \(\tau(m)=(m+1/m)/2-1\) term by term: \(\sum m/2 = N(N+1)/4\), \(\sum(1/m)/2=H_N/2\), and \(\sum 1\) over \(N\) terms is \(N\).
\end{proof}

\begin{proposition}[The wedge integral]
\label{prop:wedgeint}
\[
\int_1^N \tau(k)\,dk=\frac{N^2-1}{4}+\frac12\ln N-(N-1).
\]
\end{proposition}

\begin{proof}
\(\tau(k)=k/2-1+1/(2k)\); integrate term by term from \(1\) to \(N\).
\end{proof}

\begin{proposition}[The trapezoid residual]
\label{prop:wedge}
Write \(S(N)\) for the discrete sum of Proposition~\ref{prop:tausum}, \(I(N)\) for the integral of Proposition~\ref{prop:wedgeint}, and
\(T(N)=I(N)+\tfrac12\bigl(\tau(1)+\tau(N)\bigr)\) for the trapezoidal estimate built from it.
Let \(\gamma=\lim_{N\to\infty}(H_N-\ln N)\) be the Euler--Mascheroni constant.
Then
\[
S(N)-T(N)=\frac12\bigl(H_N-\ln N\bigr)-\frac14-\frac1{4N}
\qquad\xrightarrow[N\to\infty]{}\qquad
\frac{\gamma}{2}-\frac14 = 0.0386078\ldots
\]
\end{proposition}

\begin{proof}
\(\tau(1)=0\), so \(T(N)=I(N)+\tfrac12\tau(N)\).
Subtracting the closed forms of Propositions~\ref{prop:tausum} and~\ref{prop:wedgeint} term by term,
\[
S(N)-I(N)
=\left[\frac{N(N+1)}4+\frac{H_N}2-N\right]-\left[\frac{N^2-1}4+\frac12\ln N-(N-1)\right]
=\frac{N+1}4+\frac12\bigl(H_N-\ln N\bigr)-1.
\]
Since \(\tau(N)=(N-1)^2/(2N)=N/2-1+1/(2N)\),
\[
\frac12\tau(N)=\frac N4-\frac12+\frac1{4N}.
\]
Subtract this from \(S(N)-I(N)\):
\[
S(N)-T(N)=\left[\frac{N+1}4+\frac12(H_N-\ln N)-1\right]-\left[\frac N4-\frac12+\frac1{4N}\right]
=\frac12(H_N-\ln N)-\frac14-\frac1{4N}.
\]
Let \(N\to\infty\): \(H_N-\ln N\to\gamma\) by definition of \(\gamma\), and \(1/(4N)\to0\).
\end{proof}

A term-by-term Euler--Maclaurin expansion of \(\tau\) using only its derivatives at the two endpoints, \(\tau'(x)=\tfrac12-\tfrac1{2x^2}\), predicts a correction of \(\tfrac{B_2}2\bigl[\tau'(N)-\tau'(1)\bigr]\to\tfrac1{12}\cdot\tfrac12=\tfrac1{24}\) \cite{apostol1976}.
That reasoning is not wrong about the polynomial part of \(\tau\); it is incomplete about the part that is not local.
\(\tau\) carries a term \(1/(2x)\), and the discrete sum of \(1/m\) differs from its integral not by anything visible in the derivatives of \(1/x\) at \(1\) and \(N\), but by the Euler--Mascheroni constant --- a limit, not a local correction.
\(\gamma/2-1/4\) and \(1/24\) are close (\(0.0386\) against \(0.0417\)) and easy to mistake for one another at a single \(N\), which is exactly why the recorded run's own digits, not either asymptotic guess, are what this series has always deferred to.

On the recorded run, \(N=127\): the exact discrete sum is \(3939.712667\); the integral plus trapezoidal correction is \(3939.674062\); the difference is \(0.038605\).
Proposition~\ref{prop:wedge}'s limit, \(\gamma/2-1/4=0.0386078\ldots\), already agrees with that run to five figures at \(N=127\); \(1/24=0.041\overline6\) does not.
As stated in \emph{Before the proofs}, this sum is always taken from the discrete closed form of Proposition~\ref{prop:tausum}, never from a running total carried across many hits and left uncompared to anything --- and, now, never mistaken for the bare integral it is not.

\subsection{What every hit shares}

\begin{align*}
r\cdot(1/r)&=1,\qquad
\mathrm{GM}(r,1/r)=1,\qquad
\tau(m)=(m-1)^2/(2m),\\
\bigl(1/m,\, m;\, 1,\, -1\bigr)&=-1,\qquad
\bigl(m-2,m+2;m-1,m+1\bigr)=1/9,\\
P_m&=m\,u_m,\\
\text{opposite-wall chord (idealized)}&=2m-1,\\
\text{chord/path bands (first-order)}&=2/(\mathrm{gap}\cdot h),\\
\textstyle\sum_{m=1}^N\tau(m)&=N(N+1)/4+H_N/2-N.
\end{align*}
The alphabet does not grow.
Every later hit writes the same words at a larger scale.

% ======================================================================
\section{The other end}
\label{sec:otherend}
% ======================================================================

Section~\ref{sec:outward} built one point from the count: \(P_m=m\,u_m\), walking outward, one hit at a time.
The pair it came from had two halves.
Only one of them has been drawn.

\[
Q_m=\frac{u_m}{m}.
\]

Nothing is added to produce \(Q_m\) that was not already proved above.

\begin{proposition}
\label{prop:qisinv}
\(Q_m\) is the inversion of \(P_m\) through the bind: \(Q_m = P_m/|P_m|^2\).
\end{proposition}

\begin{proof}
\(|P_m|^2=m^2\), so \(P_m/|P_m|^2 = m\,u_m/m^2 = u_m/m = Q_m\).
\end{proof}

\(P_m\) and \(Q_m\) sit on the same ray, through the same origin, at the same instant.
They are not two walks that happen to share a clock.
They are one walk, and the classical operation of inversion through the bind, applied to it.

\begin{proposition}[Power of a point]
\label{prop:power}
\(\mathrm{pow}(P_m)\cdot\mathrm{pow}(Q_m) = -\mathrm{pow}(P_m)^2/m^2\), where \(\mathrm{pow}(X)=|OX|^2-1\) is the power of \(X\) with respect to the bind.
\end{proposition}

\begin{proof}
\(\mathrm{pow}(P_m)=m^2-1\); \(\mathrm{pow}(Q_m)=1/m^2-1=-(m^2-1)/m^2\).
Their product is \(-(m^2-1)^2/m^2=-\mathrm{pow}(P_m)^2/m^2\), an immediate corollary of Proposition~\ref{prop:qisinv}.
\end{proof}

\subsection{The inward point}

\begin{proposition}
\label{prop:pq}
\(|OP_m|\cdot|OQ_m|=1\) for every \(m\ge 1\).
\end{proposition}

\begin{proof}
\(|OP_m|=m\), \(|OQ_m|=1/m\), since \(u_m\) has length \(1\).
This is \(r\cdot(1/r)=1\), read as a distance.
\end{proof}

At \(m=1\), \(P_1=Q_1=u_1\).
For \(m>1\), \(P_m\) is further from the origin than the unit circle, and \(Q_m\) is closer, by the same walk, read in the two directions its own pair already names.

\begin{proposition}[Length reading of the leftover]
\label{prop:taulength}
\[
\tau(m)=\tfrac12\cdot\mathrm{dist}(\text{unit circle},P_m)\cdot\mathrm{dist}(\text{unit circle},Q_m).
\]
\end{proposition}

\begin{proof}
\(\mathrm{dist}(\text{unit circle},P_m)=m-1\) along the ray;
\(\mathrm{dist}(\text{unit circle},Q_m)=1-1/m=(m-1)/m\).
Their product, halved, is \((m-1)^2/(2m)\), which is \(\tau(m)\) by Proposition~\ref{prop:tau}.
\end{proof}

The leftover was defined above as a gap between two means.
It is also half the product of two lengths on the ray --- how far \(P_m\) has walked past the lock, times how far \(Q_m\) has walked back from it.
At the lock both lengths are zero and \(\tau\) is zero for the same reason twice.

\begin{observation}
\(Q_m\) requires no rule beyond the one that produced \(P_m\).
Both points exist at every hit, from the first.
Neither waits on the other to arrive.
\end{observation}

\(P_m\) grows without bound as the count runs.
\(Q_m\) shrinks toward the origin and never reaches it, since \(1/m>0\) for every finite \(m\).
The lock is the one hit at which asking the question either way gives the same answer.

\begin{figure}[htbp]
\centering
\includegraphics[width=0.92\textwidth]{S1_section.png}
\caption{Sheet S1. Section $y=0$. Outer stations on $z=|x|$; inverse stations on $z=1/|x|$, which at this scale collapse toward the count-axis. They meet at $m=1$.}
\label{fig:section}
\end{figure}

\subsection{The inward wall}

\begin{proposition}
\label{prop:innerchord}
Under the same idealization as Proposition~\ref{prop:chord} --- opposite headings at consecutive hits ---
\[
|Q_m-Q_{m-1}|=\frac{2m-1}{m(m-1)}.
\]
\end{proposition}

\begin{proof}
\[
\left|\frac{u}{m}-\frac{-u}{m-1}\right|
=\left(\frac1m+\frac1{m-1}\right)|u|
=\frac{2m-1}{m(m-1)}.
\]
\end{proof}

\begin{proposition}
\label{prop:product4}
\[
(2m-1)\cdot\frac{2m-1}{m(m-1)}
=\frac{(2m-1)^2}{m(m-1)}
\;\longrightarrow\; 4
\quad\text{as }m\to\infty.
\]
\end{proposition}

\begin{proof}
\[
\frac{(2m-1)^2}{m(m-1)}
=\frac{4m^2-4m+1}{m^2-m}
=4+\frac{1}{m^2-m}.
\]
\end{proof}

At \(m=2\) the product is \(4.5\).
On the recorded run, at \(m=127\), it is \(4.000062\).
The limit is not a value the walk equals at any finite hit.
The formula is exact algebra given the idealized chords; that those chords approximate the real, drifting ones is the entrance residual, carried forward unchanged.

\subsection{Solids from the same walk}
\label{sec:solids}

Stack \(P_m\) and \(Q_m\) by height \(z=m\), so the outward locus becomes the cone \(r=z\) and the inward locus becomes the funnel \(r=1/z\), with the bind stood up as the cylinder \(r=1\) between them, from \(z=1\) to \(z=N\).

\begin{proposition}
\label{prop:volumes}
\[
\frac{V_{\mathrm{in}}}{\pi}=\frac{(N-1)^2}{N}=2\,\tau(N),
\qquad
\frac{V_{\mathrm{out}}}{\pi}=\frac{(N-1)^2(N+2)}{3},
\qquad
\frac{V_{\mathrm{out}}}{V_{\mathrm{in}}}=\frac{N(N+2)}{3}.
\]
\end{proposition}

\begin{proof}
\(V_{\mathrm{cone}}=\tfrac{\pi}{3}(N^3-1)\), \(V_{\mathrm{cyl}}=\pi(N-1)\), \(V_{\mathrm{funnel}}=\pi(1-1/N)\).
\(V_{\mathrm{out}}=V_{\mathrm{cone}}-V_{\mathrm{cyl}}\) and \(V_{\mathrm{in}}=V_{\mathrm{cyl}}-V_{\mathrm{funnel}}\); both reduce by direct algebra to the stated forms, and \((N-1)^2/N=2\tau(N)\) by Proposition~\ref{prop:tau}.
\end{proof}

No third point is needed to reach this.
The leftover \(\tau(N)\), already read as a gap between means and as half a product of two lengths, is now also, up to a factor of \(2\pi\), the volume of the solid the inward walk alone sweeps out between the lock and hit \(N\).
Three readings of one number, none of them requiring anything the walk had not already produced.

On the recorded run, \(N=127\): \(V_{\mathrm{in}}/\pi=126^{2}/127=2\tau(127)\), and \(V_{\mathrm{out}}/V_{\mathrm{in}}=127\cdot 129/3=5461\).

\begin{figure}[htbp]
\centering
\includegraphics[width=0.92\textwidth]{O1_oblique.png}
\caption{Sheet O1. The same vertices, un-cut. Outer cone in ink; inverse funnel in the accent; bind as the count-axis.}
\label{fig:oblique}
\end{figure}

% ======================================================================
\section{Meetings}
\label{sec:meetings}
% ======================================================================

\(P_m\) and \(Q_m\) share an origin and a ray, by construction, at every hit.
What was not yet asked is whether they share anything else --- a circle, a place neither one owns alone.
This section finds two such places.
One is exact at every hit, with no idealization in it.
The other is exact only under the idealization already named, and the run confirms it to the precision that idealization allows.

\subsection{The circle on \(OP\)}
\label{sec:circleOP}

Let \(O\) be the origin and let \(C_m\) be the circle with diameter \(OP_m\): centre \(P_m/2\), radius \(m/2\).

\begin{proposition}
\label{prop:circleOP}
\(C_m\) meets the unit circle at the two points
\[
X_\pm
=\frac1m\,u_m
\;\pm\;
\sqrt{1-\frac1{m^2}}\;Ju_m,
\qquad
Ju_m=(-h_y,h_x),
\]
and the orthogonal projection of either \(X_\pm\) onto the ray of \(u_m\) is exactly \(Q_m\).
\end{proposition}

\begin{proof}
A point \(X\) lies on \(C_m\) iff \(|X-P_m/2|^2=(m/2)^2\), which expands, using \(|u_m|=1\), to \(|X|^2=X\cdot P_m\).
On the unit circle, \(|X|=1\), so the intersection satisfies \(X\cdot u_m=1/m\).
Writing \(X=(1/m)u_m+t\,Ju_m\) and imposing \(|X|=1\) forces \(t=\pm\sqrt{1-1/m^2}\).
The projection of either root onto the ray of \(u_m\) is \((X\cdot u_m)u_m=(1/m)u_m=Q_m\).
\end{proof}

Nothing was idealized to get this.
\(P_m\) is the recorded outward point at hit \(m\).
The circle it defines meets the unit circle at a point whose foot on the ray is exactly the recorded \(Q_m\).
On the recorded rows \(m=1,2,4,5,10,36,127\), the foot agrees with the declared \(Q_m\) to machine precision, and \(|X_+|=1\) to the same precision.
At \(m=1\) the two circles that could be built this way --- on \(OP_1\) and on \(OQ_1\) --- coincide, and \(X_+=X_-=P_1=Q_1=u_1\): the lock, as the one place a construction stops needing two cases.
This construction is not new: it is the standard ruler-and-compass construction of an inverse point --- construct the circle on the diameter joining the centre and the point, and its intersection with the given circle marks the inverse on the line between them \cite{coxeter1969}.
What is being read here is that construction, run on the outward point the walk itself supplies.

Thales puts a right angle at each \(X_\pm\).
The kite \(OX_+P_mX_-\) therefore has two right angles.
The other two are \(2\arccos(1/m)\) and \(2\arcsin(1/m)\).
At \(m=2\) the kite is \(90^\circ\), \(90^\circ\), \(120^\circ\), \(60^\circ\).
Scaled by the count,
\[
2m\arcsin\frac1m \;\longrightarrow\; 2.
\]
At \(m=127\) the value is \(2.000021\).
The unit wall-chord, recovered from the opposite vertex of the same two circles.

\begin{figure}[htbp]
\centering
\includegraphics[width=0.92\textwidth]{D6_kite.png}
\caption{Sheet D6. The circle on $OP_2$ meeting the unit circle at $X_\pm$; the kite $OX_+P_2X_-$ with its two right angles. The foot of either chord on the ray is $Q_2$.}
\label{fig:kite}
\end{figure}

\subsection{Inversion of a collinear diameter}
\label{sec:generallemma}

The circle-on-\(OP\) construction inverts a single point.
A second question is what inversion does to a whole circle built from two points of the walk, when those two points and the origin happen to sit on one line.

\begin{proposition}[Inversion of a collinear diameter]
\label{prop:geninv}
Let \(C\) be a circle with diameter endpoints \(A,B\), and suppose \(O\), \(A\), \(B\) are collinear.
Under inversion through a circle centred at \(O\), the image of \(C\) is the circle with diameter endpoints \(\mathrm{Inv}(A)\) and \(\mathrm{Inv}(B)\).
\end{proposition}

\begin{proof}
The line through \(O,A,B\) is an axis of symmetry of \(C\), since any diameter is.
Inversion centred at \(O\) commutes with reflection across any line through \(O\): reflecting first and inverting after gives the same point as inverting first and reflecting after, since both act on each ray through \(O\) independently.
So \(\mathrm{Inv}(C)\) inherits that same axis of symmetry, which forces its centre onto that line.
\(\mathrm{Inv}(C)\) also passes through \(\mathrm{Inv}(A)\) and \(\mathrm{Inv}(B)\), themselves on that line.
A circle meets any line through its own centre in exactly two points, and those two points are automatically antipodal.
So \(\mathrm{Inv}(A)\) and \(\mathrm{Inv}(B)\) are forced to be \(\mathrm{Inv}(C)\)'s diameter endpoints.
\end{proof}

Two hits \(m,n\) supply such a collinear pair whenever their headings are, to the entrance residual, parallel or antiparallel --- same wall or opposite wall, at any separation, not only adjacent hits.

\begin{proposition}[The swap, in general]
\label{prop:swap}
Under the idealization that \(u_m\) and \(u_n\) are equal or opposite, inversion through the unit circle sends the diameter-circle of \(P_m,P_n\) onto the diameter-circle of \(Q_m,Q_n\).
\end{proposition}

\begin{proof}
Immediate from Proposition~\ref{prop:geninv} with \(A=P_m\), \(B=P_n\), together with Proposition~\ref{prop:qisinv} identifying \(\mathrm{Inv}(P_m)=Q_m\) and \(\mathrm{Inv}(P_n)=Q_n\).
\end{proof}

Checked against the recorded run at same-wall pairs \((m,n)=(2,4),(5,13),(5,25),(10,50),(3,127)\) and opposite-wall pairs \((5,14),(5,26),(10,51),(3,126)\), predicted and actual inner centres and radii agree to between four and twelve significant figures in every case.
The precision is not uniform, and it is not supposed to be: it is governed entirely by how close \(u_m\) and \(u_n\) are to exactly parallel or antiparallel at that pair, which is the same ratchet-and-reset drift of \S\ref{sec:ratchet}, not a separate source of error.
Proposition~\ref{prop:swap} is not a fact about consecutive same-wall hits.
It is a fact about any two collinear points the walk produces, of which consecutive same-wall and consecutive opposite-wall hits are the two smallest and most easily checked cases.

\begin{proposition}[Clearance from the bind]
\label{prop:clearance}
For the same-wall pair \(m,m+2\), under the idealization above, the outer diameter-circle stands off the unit circle by exactly \(m-1\) on its near side and \(m+1\) on its far side.
The inner diameter-circle stands off it by exactly \(1-1/m\) on its near side and \(1-1/(m+2)\) on its far side.
\end{proposition}

\begin{proof}
The outer circle has centre at distance \(m+1\) and radius \(1\); its near and far points from the origin sit at \(m+1-1=m\) and \(m+1+1=m+2\), so their distances from the unit circle are \(m-1\) and \(m+1\).
The inner circle has centre at distance \((m+1)/(m(m+2))\) and radius \(1/(m(m+2))\); its near and far points sit at \(1/(m+2)\) and \(1/m\), so their distances from the unit circle are \(1-1/(m+2)\) and \(1-1/m\).
\end{proof}

The diameter-circle of \(P_mP_{m+2}\) does not meet the unit circle for \(m\ge 2\).
The diameter-circle of \(Q_mQ_{m+2}\) does not meet it for \(m>1\).
At \(m=1\) both touch at the lock, and Proposition~\ref{prop:clearance}'s clearances both fall to zero, consistently.

\begin{figure}[htbp]
\centering
\includegraphics[width=0.92\textwidth]{D7_swap.png}
\caption{Sheet D7. Inversion swapping the outer diameter-circle of the same-wall pair $P_2,P_4$ onto the inner diameter-circle of $Q_2,Q_4$.}
\label{fig:swap}
\end{figure}

\subsection{What does not meet}
\label{sec:notmeet}

\(P_m\) grows without bound as \(m\) runs.
\(Q_m\) shrinks toward the origin and never reaches it.
The two readings share every hit and share no destination.

The chord product of the inward and outward walls approaches \(4\) and is not equal to it at any finite \(m\).
A meeting that is forced holds at every hit without qualification, as the circle on \(OP\) does.
A meeting that depends on an idealization holds only as well as the idealization does, as the swap of Proposition~\ref{prop:swap} does --- and that idealization's accuracy varies with which pair is chosen, not with how far apart they are.
Where the two readings end up together is not a question with a yes, because one of them does not end up anywhere at all.

A smaller mismatch belongs here too, because it is the same discipline at a different scale.
Stand the same-wall span \(|P_{m+2}-P_m|\) up by height \(z=m\), so that \(P_m\) and \(P_{m+2}\) become \((m,m)\) and \((m+2,m+2)\) in a plane cut through the count.
Looking down that count --- the plan --- the span between them is \(2\).
Cut open to that plane instead --- the section --- the same two points sit on a line of slope \(1\) and are \(2\sqrt2\) apart, not \(2\).
Both numbers are exact; neither is a rounding of the other.
A quantity can be genuinely invariant in the object and still read differently depending on which projection of it is taken, and the two readings above are not in tension for the same reason the plan and the section of any solid are not in tension: they are answers to different questions asked of the same figure.

% ======================================================================
\section{What stands}
% ======================================================================

Forced, at every hit, once the operators of \S1 are granted:
the shear \(T_h\) and the bind \(N\), with no angle in either;
a hit as a sign change on the diameter the heading stood on;
the pair \(r=m\), \(1/r=1/m\), product \(1\), geometric mean \(1\);
the leftover \(\tau(m)=(m-1)^2/(2m)\), zero only at \(m=1\), and its closed-form sum;
the point \(P_m=m\,u_m\) and the point \(Q_m=u_m/m\);
\(Q_m\) as the inversion of \(P_m\) through the bind, and every fact below as a consequence of that one classical operation applied to the outward walk;
their product \(|OP_m||OQ_m|=1\) and their power-of-a-point product;
the leftover, again, as half a product of two lengths, and again as a volume;
four points on the cut writing \(-1\);
the generator ladder writing \(1/9\);
the circle on \(OP_m\) meeting the unit circle at a point whose foot is \(Q_m\) --- the classical construction of an inverse point, run on the walk's own output, with no idealization, at every hit including the lock;
the general lemma that inversion of a circle whose diameter is collinear with the centre of inversion sends it to the circle on the images of that diameter's endpoints, forced for any two such points, not only ones the walk happens to produce.

Forced under the idealization that two hits share a parallel or antiparallel heading, and confirmed by the run at whatever precision that idealization allows for the pair chosen:
the outward and inward chords, their product converging to \(4\);
the swap of Proposition~\ref{prop:swap}, checked at separations from \(2\) to \(124\);
the exact clearances of Proposition~\ref{prop:clearance}.

Reported from the table, and compared after the table exists:
two gap integers, with a single-letter alphabet at the exact mixes of Observation~\ref{obs:exactclosure};
mean gap times kick sitting next to \(\pi\);
two first-order chord/path bands with \(2/\pi\) between them;
\(\alpha\) on short hits of this run;
the plan-and-section mismatch of the same-wall span.

Named: entrance residual, ratchet and reset, the inward point, the general lemma, the foot of the chord, the generator ladder.

Not granted:
an angle as an argument of either move;
an exact opposite heading at consecutive hits, only an approach to one;
a value of \(\pi\) or of \(4\) treated as reached rather than approached;
a second move, or any rule governing \(Q_m\) beyond the one that produced \(P_m\);
a meeting of the two readings' destinations;
a third point, where two already account for everything checked;
a running sum trusted over its own closed form.

The first sentence of the series is still the sentence.
Where is \(1\)?
It is the lock: the one hit where the pair meets itself, where \(P\) and \(Q\) are still the same point, and where the two circles of \S\ref{sec:meetings} are still one circle --- before the count has given any of them somewhere else to be.

\newpage
\begin{thebibliography}{9}

\bibitem{coxeter1969}
H.\,S.\,M.\ Coxeter, \emph{Introduction to Geometry}, 2nd ed., Wiley, 1969.

\bibitem{sos1958}
V.\,T.\ S\'os, ``On the theory of diophantine approximations I,'' \emph{Acta Mathematica Academiae Scientiarum Hungaricae} 8 (1958), 461--472.

\bibitem{swierczkowski1958}
S.\ \'Swierczkowski, ``On successive settings of an arc on the circumference of a circle,'' \emph{Fundamenta Mathematicae} 46 (1958), 187--189.

\bibitem{volder1959}
J.\,E.\ Volder, ``The CORDIC Trigonometric Computing Technique,'' \emph{IRE Transactions on Electronic Computers} EC-8, no.\,3 (1959), 330--334.

\bibitem{apostol1976}
T.\,M.\ Apostol, \emph{Introduction to Analytic Number Theory}, Springer, 1976.

\end{thebibliography}

\end{document}
